\section{Piecewise quadratic supermartingales in multiple dimensions}

Now we will finally consider a very general setup. 

First, we prove that when a piecewise-quadratic function $V : \mathbb{R}^d \to \mathbb{R}$ 
satisfies similar concativity and smoothness conditions as in the one-dimensional case, then $V(X_t)$ is a supermartingale.

Second, we will prove that if a smooth certificate exists, then we can construct a piecewise quadratic certificate that is a supermartingale,
essentially showing how quadratic certificates are in some sense sufficient, while piecewise linear certificates are not.

\subsection{Proving the supermartingale property}

We will show that a piecewise quadratic function gives a supermartingale if:
\begin{enumerate}
\item the generator is nonpositive inside each region,
\item the stochastic integral is a martingale,
\item the normal derivative jumps in the correct direction across each
interface.
\end{enumerate}

\begin{setup}
Let:
\begin{enumerate}
\item $dX_t = f(X_t)dt + \sigma(X_t)dW_t$ be a $d$-dimensional SDE, where
$f:\mathbb{R}^d \to \mathbb{R}^d$ and
$\sigma:\mathbb{R}^d \to \mathbb{R}^{d \times m}$.
\item $W_t$ is an $m$-dimensional Brownian motion.
\item $a(x) = \sigma(x)\sigma(x)^\top$.
\item $K \subseteq \mathbb{R}^d$ is a compact set partitioned into finitely
many polyhedral cells:
\[
K = \bigcup_{i=1}^N K_i.
\]
\item $V:K \to \mathbb{R}$ is a continuous piecewise quadratic function such
that
\[
V(x) = V_i(x)
=
c_i + p_i^\top x + \frac{1}{2}x^\top Q_i x,
\qquad x \in K_i,
\]
where $Q_i = Q_i^\top$.
\item $V$ is continuous across interfaces:
\[
V_i(x) = V_j(x),
\qquad x \in K_i \cap K_j.
\]
\end{enumerate}
\end{setup}

Now we will consider what conditions on $V$ are sufficient to guarantee that
$V(X_t)$ is a supermartingale, at least up to the first exit time from $K$.

Define the stopping time
\[
\tau_K = \inf\{t \geq 0 : X_t \notin K\}.
\]

\begin{hypothesis}
Does a multidimensional analogue of concavity, together with an interior
generator condition, guarantee that
\[
V(X_{t \wedge \tau_K})
\]
is a supermartingale?
\end{hypothesis}

To answer this, we apply the multidimensional Itô--Tanaka formula to
$V(X_{t \wedge \tau_K})$.

On each cell $K_i$, we have
\[
\nabla V_i(x) = p_i + Q_i x,
\qquad
\nabla^2 V_i(x) = Q_i.
\]

Therefore, the infinitesimal generator of $V_i$ is
\[
\mathcal{L}V_i(x)
=
\nabla V_i(x) \cdot f(x)
+
\frac{1}{2}
\operatorname{tr}\!\left(a(x)\nabla^2 V_i(x)\right).
\]

Substituting the quadratic form gives
\[
\mathcal{L}V_i(x)
=
(p_i + Q_i x)\cdot f(x)
+
\frac{1}{2}\operatorname{tr}\!\left(a(x)Q_i\right).
\]

Now let $F_{ij} = K_i \cap K_j$ be a shared interface between two neighbouring
cells. Let $n_{ij}$ be the unit normal vector pointing from $K_i$ into $K_j$.

The multidimensional Itô--Tanaka formula gives
\begin{align}
V(X_{t \wedge \tau_K})
=
V(X_0)
&+
\int_0^{t \wedge \tau_K}
\mathcal{L}V(X_s) ds \\
&+
\int_0^{t \wedge \tau_K}
\nabla V(X_s)^\top \sigma(X_s) dW_s \\
&+
\frac{1}{2}
\sum_{i<j}
\int_0^{t \wedge \tau_K}
\left(
\nabla V_j(X_s) - \nabla V_i(X_s)
\right)\cdot n_{ij}
dL_s^{F_{ij}} .
\end{align}

Here $L_t^{F_{ij}}$ denotes the surface local time accumulated by $X_t$ on the
interface $F_{ij}$.

We will focus on these terms:

\begin{enumerate}
\item
\[
D_t
=
\int_0^{t \wedge \tau_K}
\mathcal{L}V(X_s)ds,
\]
the interior drift term.

\item
\[
    M_t
    =
    \int_0^{t \wedge \tau_K}
    \nabla V(X_s)^\top \sigma(X_s)dW_s,
\]
the diffusion term.

\item
\[
    K_t
    =
    \frac{1}{2}
    \sum_{i<j}
    \int_0^{t \wedge \tau_K}
    \left(
        \nabla V_j(X_s) - \nabla V_i(X_s)
    \right)\cdot n_{ij}
    dL_s^{F_{ij}},
\]
the interface term.
\end{enumerate}

We need to prove that
\[
D_t + M_t + K_t
\]
is a supermartingale.

Some natural assumptions, which are sufficient to guarantee that
$V(X_{t \wedge \tau_K})$ is a supermartingale, are the following.

\begin{enumerate}
\item
The interior generator is nonpositive on every cell:
\[
\mathcal{L}V_i(x) \leq 0,
\qquad x \in \operatorname{int}(K_i).
\]
Equivalently,
\[
(p_i + Q_i x)\cdot f(x)
+
\frac{1}{2}\operatorname{tr}\!\left(a(x)Q_i\right)
\leq 0,
\qquad x \in \operatorname{int}(K_i).
\]

This is the multidimensional analogue of the condition
\[
    f(x)V'(x) \leq 0
\]
from the one-dimensional piecewise linear case.

\begin{fact}
If $\mathcal{L}V(X_s) \leq 0$ for all $s \leq \tau_K$, then the drift term
$D_t$ is a supermartingale.

\textbf{Proof:} We want to prove
\[
    \mathbb{E}[D_t \mid \mathcal{F}_r] \leq D_r,
    \qquad r \leq t.
\]
By definition,
\begin{align}
    \mathbb{E}[D_t \mid \mathcal{F}_r]
    &=
    \mathbb{E}\left[
    \int_0^{t \wedge \tau_K}
    \mathcal{L}V(X_s)ds
    \Bigm| \mathcal{F}_r
    \right] \\
    &=
    \int_0^{r \wedge \tau_K}
    \mathcal{L}V(X_s)ds
    +
    \mathbb{E}\left[
    \int_{r \wedge \tau_K}^{t \wedge \tau_K}
    \mathcal{L}V(X_s)ds
    \Bigm| \mathcal{F}_r
    \right].
\end{align}
Since $\mathcal{L}V(X_s) \leq 0$, the second term is nonpositive. Hence
\[
    \mathbb{E}[D_t \mid \mathcal{F}_r]
    \leq
    D_r.
\]
Therefore $D_t$ is a supermartingale.
\end{fact}

\item
The stochastic integral $M_t$ is a martingale. A sufficient condition is
\[
    \mathbb{E}\left[
    \int_0^{t \wedge \tau_K}
    \left\|
        \nabla V(X_s)^\top \sigma(X_s)
    \right\|^2 ds
    \right]
    < \infty.
\]
Since $V$ is piecewise quadratic on finitely many cells, $\nabla V$ is
piecewise linear. Therefore, on a compact set $K$, $\nabla V$ is bounded.
Thus this condition is usually easy to verify if $\sigma$ is bounded on $K$.

\item
The interface contribution is nonpositive. For every shared face
\[
    F_{ij} = K_i \cap K_j,
\]
we assume
\[
    \left(
        \nabla V_j(x) - \nabla V_i(x)
    \right)\cdot n_{ij}
    \leq 0,
    \qquad x \in F_{ij}.
\]
Since
\[
    \nabla V_i(x) = p_i + Q_i x
\]
and
\[
    \nabla V_j(x) = p_j + Q_j x,
\]
this condition can be written explicitly as
\[
    \left(
        p_j - p_i + (Q_j - Q_i)x
    \right)\cdot n_{ij}
    \leq 0,
    \qquad x \in F_{ij}.
\]

This is the multidimensional version of the concavity condition
\[
    V'_+(x_i) - V'_-(x_i) \leq 0
\]
from the one-dimensional case.

\begin{fact}
If
\[
    \left(
        \nabla V_j(x) - \nabla V_i(x)
    \right)\cdot n_{ij}
    \leq 0
    \qquad \text{on } F_{ij},
\]
and $L_t^{F_{ij}}$ is nondecreasing in $t$, then the interface term $K_t$ is
a supermartingale.

\textbf{Proof:} We want to prove
\[
    \mathbb{E}[K_t \mid \mathcal{F}_r] \leq K_r,
    \qquad r \leq t.
\]
The interface term is
\[
    K_t
    =
    \frac{1}{2}
    \sum_{i<j}
    \int_0^{t \wedge \tau_K}
    \gamma_{ij}(X_s)dL_s^{F_{ij}},
\]
where
\[
    \gamma_{ij}(x)
    =
    \left(
        \nabla V_j(x) - \nabla V_i(x)
    \right)\cdot n_{ij}.
\]
By assumption,
\[
    \gamma_{ij}(x) \leq 0
    \qquad x \in F_{ij}.
\]
Also, $dL_s^{F_{ij}}$ is a nonnegative measure, because surface local time is
nondecreasing. Therefore
\[
    \int_r^t \gamma_{ij}(X_s)dL_s^{F_{ij}} \leq 0.
\]
Hence the future increment of $K_t$ after time $r$ is nonpositive, and so
\[
    \mathbb{E}[K_t \mid \mathcal{F}_r] \leq K_r.
\]
Therefore $K_t$ is a supermartingale.
\end{fact}
\end{enumerate}

\subsubsection{Proof of the supermartingale property}

We now combine the three conditions.

By the multidimensional Itô--Tanaka formula,
\begin{align}
V(X_{t \wedge \tau_K})
=
V(X_0)
+
D_t
+
M_t
+
K_t.
\end{align}

The drift term $D_t$ is a supermartingale by the interior generator condition:
\[
\mathcal{L}V_i(x) \leq 0
\qquad x \in \operatorname{int}(K_i).
\]

The stochastic integral $M_t$ is a martingale by the square-integrability
assumption:
\[
\mathbb{E}\left[
\int_0^{t \wedge \tau_K}
\left|
\nabla V(X_s)^\top \sigma(X_s)
\right|^2 ds
\right]
< \infty.
\]

The interface term $K_t$ is a supermartingale by the multidimensional concavity
condition:
\[
\left(
\nabla V_j(x) - \nabla V_i(x)
\right)\cdot n_{ij}
\leq 0,
\qquad x \in F_{ij}.
\]

Therefore $D_t + K_t$ is a supermartingale, and $M_t$ is a martingale. Since the
sum of a martingale and a supermartingale is a supermartingale, we conclude that
\[
V(X_{t \wedge \tau_K})
\]
is a supermartingale.

\subsubsection{Interpretation of the multidimensional concavity condition}

The condition
\[
\left(
\nabla V_j(x) - \nabla V_i(x)
\right)\cdot n_{ij}
\leq 0,
\qquad x \in F_{ij},
\]
says that when the process crosses from cell $K_i$ into cell $K_j$, the
directional derivative of $V$ in the crossing direction cannot increase.

Indeed, the normal derivative from the $K_i$ side is
\[
\nabla V_i(x)\cdot n_{ij},
\]
while the normal derivative from the $K_j$ side is
\[
\nabla V_j(x)\cdot n_{ij}.
\]

Thus the condition is equivalent to
\[
\nabla V_j(x)\cdot n_{ij}
\leq
\nabla V_i(x)\cdot n_{ij}.
\]

This is exactly the analogue of the one-dimensional condition
\[
V'_+(x_i) \leq V'_-(x_i).
\]

So, in multiple dimensions, ordinary concavity is not the only thing we need to
look at. What matters for the local time term is the jump in the normal
derivative across each interface. If that jump is nonpositive, then the surface
local time contribution pushes $V(X_t)$ downwards rather than upwards.

\subsubsection{Summary of sufficient conditions}

We have shown that the following assumptions are sufficient.

\begin{enumerate}
\item $V$ is continuous and piecewise quadratic:
\[
V(x)
=
c_i + p_i^\top x + \frac{1}{2}x^\top Q_i x,
\qquad x \in K_i.
\]

\item The generator is nonpositive inside each cell:
\[
    (p_i + Q_i x)\cdot f(x)
    +
    \frac{1}{2}\operatorname{tr}\!\left(a(x)Q_i\right)
    \leq 0,
    \qquad x \in \operatorname{int}(K_i).
\]

\item The stochastic integral is a martingale:
\[
    \mathbb{E}\left[
    \int_0^{t \wedge \tau_K}
    \left\|
        \nabla V(X_s)^\top \sigma(X_s)
    \right\|^2 ds
    \right]
    < \infty.
\]

\item The interface condition holds on every shared face:
\[
    \left(
        p_j - p_i + (Q_j - Q_i)x
    \right)\cdot n_{ij}
    \leq 0,
    \qquad x \in F_{ij}.
\]
\end{enumerate}

Under these assumptions,
\[
V(X_{t \wedge \tau_K})
\]
is a supermartingale.

This gives the multidimensional version of the one-dimensional piecewise linear
certificate result. The interior generator condition controls the ordinary Itô
drift, while the interface concavity condition controls the Tanaka local time
terms.

\subsection{Asymptotic completeness of piecewise quadratic certificates}

\paragraph{Theorem (asymptotic completeness under a strict smooth certificate).}
Let \(K\subset\mathbb R^2\) be compact and polygonal, and consider the stopped
SDE
\[
  dX_t=f(X_t)\,dt+\sigma(X_t)\,dW_t,
  \qquad
  a(x)=\sigma(x)\sigma(x)^\top,
\]
stopped at
\[
  \tau_K=\inf\{t\ge0:X_t\notin K\}.
\]
Assume that \(f\) and \(a\) are bounded on \(K\).  Suppose there exists a
strict smooth certificate \(U\in C^3(K)\) such that
\[
  \mathcal LU(x)
  =
  \nabla U(x)\cdot f(x)
  +
  \frac12\operatorname{tr}\!\bigl(a(x)\nabla^2U(x)\bigr)
  \le -\varepsilon
  \qquad x\in K,
\]
for some \(\varepsilon>0\).  Assume also that the boundary/separation
conditions hold with margin: for some \(\delta>0\),
\[
  \sup_{x\in X_0}U(x)\le \alpha-\delta,
  \qquad
  \inf_{x\in X_u}U(x)\ge \beta+\delta,
  \qquad
  \alpha<\beta.
\]

Let \(\{\mathcal T_h\}_{h\downarrow0}\) be a shape-regular sequence of
refined triangulations of \(K\) supporting \(C^1\) piecewise-quadratic spline
interpolants \(I_hU\) satisfying
\[
  \|I_hU-U\|_{C^2(K)}\to0
  \qquad\text{as }h\to0.
\]
Then, for all sufficiently small \(h\), the function
\[
  V_h:=I_hU
\]
is a continuous \(C^1\) piecewise-quadratic certificate satisfying
\[
  \mathcal LV_h(x)\le -\frac{\varepsilon}{2}
  \qquad x\in K,
\]
and
\[
  \sup_{x\in X_0}V_h(x)\le \alpha,
  \qquad
  \inf_{x\in X_u}V_h(x)\ge \beta.
\]
Consequently,
\[
  V_h(X_{t\wedge\tau_K})
\]
is a supermartingale.

\paragraph{Proof.}
Since \(f\) and \(a\) are bounded on \(K\), define
\[
  F:=\sup_{x\in K}\|f(x)\|,
  \qquad
  A:=\sup_{x\in K}\|a(x)\|.
\]
Let
\[
  E_h:=V_h-U.
\]
Then
\[
  \mathcal LV_h-\mathcal LU
  =
  \nabla E_h\cdot f
  +
  \frac12\operatorname{tr}\!\bigl(a\nabla^2E_h\bigr).
\]
Hence, for a constant \(C_K>0\) depending only on the dimension and the chosen
matrix norm,
\[
  |\mathcal LV_h(x)-\mathcal LU(x)|
  \le
  F\|\nabla E_h\|_{L^\infty(K)}
  +
  C_K A\|\nabla^2E_h\|_{L^\infty(K)}.
\]
Because
\[
  \|E_h\|_{C^2(K)}\to0,
\]
there exists \(h_0>0\) such that, for all \(h<h_0\),
\[
  |\mathcal LV_h(x)-\mathcal LU(x)|
  \le \frac{\varepsilon}{2}
  \qquad x\in K.
\]
Since \(\mathcal LU\le-\varepsilon\), it follows that
\[
  \mathcal LV_h
  \le
  -\varepsilon+\frac{\varepsilon}{2}
  =
  -\frac{\varepsilon}{2}.
\]

Similarly, since \(V_h\to U\) uniformly, for all sufficiently small \(h\),
\[
  \|V_h-U\|_{L^\infty(K)}\le\delta.
\]
Therefore
\[
  \sup_{X_0}V_h
  \le
  \sup_{X_0}U+\delta
  \le
  \alpha,
\]
and
\[
  \inf_{X_u}V_h
  \ge
  \inf_{X_u}U-\delta
  \ge
  \beta.
\]

It remains to check the nonsmooth interfaces.  The spline \(V_h\) is
\(C^1\) across every internal face of the triangulation.  Therefore, if two
quadratic pieces \(V_i,V_j\) meet across a face \(F_{ij}\), then
\[
  \nabla V_i(x)=\nabla V_j(x)
  \qquad x\in F_{ij}.
\]
Hence the Tanaka normal-derivative jump is
\[
  \bigl[\nabla V_j(x)-\nabla V_i(x)\bigr]\cdot n_{i\to j}=0,
\]
so every local-time contribution is zero.

Applying the multidimensional It\^o--Tanaka formula gives
\[
  V_h(X_{t\wedge\tau_K})
  =
  V_h(X_0)
  +
  \int_0^{t\wedge\tau_K}\mathcal LV_h(X_s)\,ds
  +
  \int_0^{t\wedge\tau_K}\nabla V_h(X_s)\sigma(X_s)\,dW_s.
\]
The drift term is nonpositive because
\[
  \mathcal LV_h\le-\frac{\varepsilon}{2}<0.
\]
The stochastic integral is a true martingale after stopping, since \(K\) is
compact and \(V_h,\sigma\) are bounded on \(K\).  Therefore
\[
  V_h(X_{t\wedge\tau_K})
\]
is a martingale plus a nonincreasing finite-variation process, and is hence a
supermartingale. \(\square\)
